\scanpage{4}
$X$ de type fini sur $Y$. On dit que $X$ est \unsure{monique|morique|monogène} sur $Y$ si

(i) les fibres de $X$ sur les $y\in Y$ sont \deleted{\unsure{toutes}} \unsure{ult\textsuperscript{t}|ab\textsuperscript{t}|rel\textsuperscript{t}} irréductibles

\inserted{alors la fonction $y\to\varphi(y)={}$\unsure{p\textsuperscript{t} gén.|p\textsuperscript{t} génér.} de $f^{-1}(y)$, est continue ssi $X\to Y$ est ouvert}

(ii) la fonction $\varphi$ est \unsure{la \uline{même}|\uline{bornée}}.

Une application $\varphi:Y\to X$ est dite \unsure{monisante|monisente|monogène} s'il existe un ouvert $U$ de $X$ qui est \inserted{\unsure{universell\textsuperscript{mt}|universellement}} monique sur $Y$,
\inserted{(i.e. $U$ \unsure{univ\textsuperscript{t}|universellement} ouvert sur $Y$)} tel que pour \deleted{\unsure{tout}} $y\in Y$,
$\varphi(y)$ soit le point générique de $U\cap f^{-1}(y)$. [Si $Y$ est noethérien, il revient au $=$ de dire : (i) Pour \deleted{\unsure{tout}} $y\in Y$,
$\varphi(y)$ est le p\textsuperscript{t} générique d'une composante \unsure{relativem\textsuperscript{t}|rel\textsuperscript{t}} irréductible de $f^{-1}(y)$ (ii) Si $y\in\overline{y'}$,
\deleted{\unsure{alors $\varphi(y')$ est l'unique} et \unsure{si} $x'\in f^{-1}(y')$ est \unsure{un point de} $f^{-1}(\unsure{y'|y})$ qui \unsure{générise} $\varphi(y)$ (iii) $\varphi$ \unsure{continue} tel que $\varphi(y)\in\overline{x'}$, alors $x'=\varphi(y')$.} \deleted{\unsure{iii}.} $\unsure{U|u}\to Y$ est \unsure{univ\textsuperscript{t}} ouvert.]
% slanted
\margin{\uline{Lemme} $\varphi$ \unsure{satisfaisant} (i) et (ii) implique l'existence de $U$ satisfaisant \unsure{(i)|(ii)}}
% slanted
\margin{C'est une propriété \uline{locale}.}

Si $\varphi:Y\to X$ est monisante, alors pour tout changement de base $Y'\to Y$, l'application associée $\varphi':Y'\to X'$ [définie de façon évidente] est monisante. Ainsi, l'ens
$F(X'/Y')$ des sections monisantes de $X'/Y'$ est un foncteur contravariant en $Y'$. Est-il représentable ??

Si oui, le \unsure{problème|préschéma} des modules $M$ sur $Y$ est nécessairement étale de type fini sur $Y$ [mais pas néc. séparé sur $Y$]

Soit $M$ \uline{étale} de type fini sur $Y$, et $\varphi_M$ une \deleted{section} application monisante pour $X\times_Y M\,/\,M$, qui ait la propriété suivante : Pour \deleted{tt} $y\in Y$ et une extension alg. close $\Omega$ de $k(y)$, l'application $M(\Omega)\to{}$ ens des composantes irréductibles de $X\otimes_y\Omega$ est \uline{bijective}. Alors $(M,\varphi_M)$ représente le foncteur $F$.
% slanted
\margin{\uline{\unsure{Car|Cor}} les \deleted{\unsure{pts}} \inserted{\uline{\unsure{éléments|étl\textsuperscript{mts}}}} de $M$ correspondent aux composantes irréductibles des fibres de $X\to Y$}
% archive sheet number: 3
